\chapter{Conflict-Based Search}
\label{ch:ch09}
% Function names for pseudocode are prefixed with "Cbs" to stay unique
% across the book.
\SetKwFunction{CbsHighLevel}{CBS}
\SetKwFunction{CbsLowLevel}{CbsLowLevel}
\SetKwFunction{CbsFirstConflict}{FirstConflict}
\SetKwFunction{CbsSplit}{Split}

\begin{objectives}
\begin{itemize}
  \item Explain the two-level idea of conflict-based search: a high-level best-first search over a binary \emph{constraint tree} and a low-level single-agent space-time \astar under constraints, and say why this beats search in the joint state space for loosely coupled teams.
  \item Define constraints, constraint-tree nodes, conflicts and the split rule exactly, and trace \cbs by hand on a small instance, drawing the constraint tree with its costs.
  \item Prove that \cbs is complete on solvable instances and returns a plan with the minimum sum of costs, and explain what its running time depends on.
  \item Implement \cbs from scratch in \python, certify its optimality against a brute-force search on tiny instances, and test it on 2, 4 and 8 agents.
  \item Recognise the classic implementation mistakes, and know the improvements (cardinal conflicts, bypass, high-level heuristics, symmetry breaking) that turn textbook \cbs into a practical solver.
\end{itemize}
\end{objectives}

% ---------------------------------------------------------------------
\section{Why this matters for a drone swarm}
\label{sec:ch09-motivation}

Twelve delivery drones share the airspace over a city block. Their routes
have been reduced, as in \cref{ch:ch02}, to a grid of cells at a fixed
altitude, with buildings as obstacles and a time step of a few seconds.
Each drone alone is easy to route: \astar from \cref{ch:ch04} finds its
shortest path in milliseconds. The difficulty is that the twelve shortest
paths cross. Two drones want the same crossing cell at the same time; a
third flies up a street while a fourth flies down it. The multi-agent path
finding (MAPF) problem of \cref{ch:ch07} asks for twelve paths without any
vertex or edge conflict, and, because every second of flight costs battery,
for paths whose \emph{sum of costs} is as small as possible.

\Cref{ch:ch07} showed that the obvious exact method, \astar in the joint
state space, does not scale: with five actions per drone (four moves and a
wait), a joint state has $5^{12} \approx 2.4 \times 10^{8}$ successors, and
the number of joint states is $\abs{V}^{12}$. \Cref{ch:ch08} showed the
opposite trade-off: prioritized planning plans the drones one after another
against a reservation table, is very fast, but can fail on solvable instances
and returns no guarantee on the cost of its plans. Conflict-based search
(\cbs)\index{conflict-based search}\index{CBS} by Sharon, Stern, Felner and
Sturtevant~\cite{sharon2012cbs,sharon2015cbs} sits between these two
extremes. It plans every drone alone, exactly as prioritized planning does,
but when two paths conflict it does not pick a winner: it tries \emph{both}
ways of resolving the conflict, by forbidding one drone or the other from the
contested cell at the contested time. It keeps both alternatives and
always continues with the cheapest world it has ever created. The result is a plan that is provably optimal and a
search that pays only for the conflicts it actually meets, not for the number
of drones. \cbs and its bounded-suboptimal sibling \ecbs (\cref{ch:ch10}) are
the standard optimal and near-optimal MAPF solvers today and the basis of most
of the research literature since 2012.

In the four-layer hybrid architecture of \cref{ch:ch24}, \cbs is the
\emph{global swarm planner}: before take-off it computes the nominal
time-indexed route of every drone, and these routes are what the prediction,
local-avoidance and replanning layers later protect and repair. Whenever a
drone has been pushed off its route by an intruder and replans locally, the
architecture re-validates the new route against the routes of the other
drones with exactly the conflict check that \cbs uses at every node. This
chapter is Week~4 of the study plan (\cref{ch:appA}): its coding milestone is
a working optimal \cbs solver, written from scratch and tested on two, four
and eight agents.

The chapter proceeds as follows. \Cref{sec:ch09-intuition} gives the
two-level idea, \cref{sec:ch09-problem} the exact definitions of constraints
and constraint-tree nodes, and \cref{sec:ch09-algorithm} the pseudocode of
both levels with a walkthrough. \Cref{sec:ch09-example} runs \cbs by hand on
a crossing of two flight corridors and draws the whole constraint tree.
\Cref{sec:ch09-properties} proves completeness and optimality and discusses
what the running time depends on. \Cref{sec:ch09-variants} presents the
improvements that make \cbs practical and an experiment on random grids,
\cref{sec:ch09-implementation} the \python implementation and its pitfalls,
and \cref{sec:ch09-drone} the place of \cbs in the drone system.

% ---------------------------------------------------------------------
\section{The idea in plain words}
\label{sec:ch09-intuition}

Imagine the drones planned in separate rooms, each unaware of the others.
Every drone comes back with its own shortest route. Now lay the routes on one
map and look for the first moment at which two drones, say $a_1$ and $a_2$,
occupy the same cell $v$ at the same time $t$. Somebody has to give way, and
we do not know who should. \cbs refuses to guess. It creates two
\emph{worlds}: in the first world a single new rule says ``$a_1$ may not be
at $v$ at time $t$'', in the second world the rule says the same about $a_2$.
In each world only the drone that received the rule goes back to its room and
replans; everybody else keeps their route. A rule can only make a route
longer or leave it unchanged, so each world is at least as expensive as its
parent.

\cbs keeps all worlds it has created in a priority queue, ordered by the sum
of the route costs, and always examines the cheapest one. If the cheapest
world has no conflict, its routes are the answer, and they are optimal: every
other world costs at least as much, and adding rules never makes a world
cheaper. If the cheapest world still has a conflict, it is split again, the
two new worlds inherit all rules of their parent plus one, and the queue is
consulted again. The set of worlds forms a binary tree, the
\textbf{constraint tree}\index{constraint tree} (CT), whose root is the world
without rules. \Cref{fig:ch09-idea} shows the first split.

\begin{figure}[tb]
  \centering
  \inputfigure{ch09/idea}
  \caption{The two-level idea. Left: planned alone, two shortest paths meet
    in cell $v$ at time $t$. Right: the high level splits the conflict into
    two children; each child adds one constraint and replans only the
    constrained agent with the low level. The high level then continues with
    the cheapest node of the whole queue, which may be either child or a
    node generated elsewhere in the tree.}
  \label{fig:ch09-idea}
\end{figure}

\begin{keyidea}
\cbs separates \emph{who yields} from \emph{how to yield}. The high level
decides who yields by searching a binary tree of constraint sets best-first
on the sum of costs; the low level decides how to yield by running
single-agent space-time \astar under the constraints of one agent. Every
split keeps at least one optimal solution, and costs never decrease down the
tree, so the first conflict-free node popped is optimal.\index{CBS!key idea}
\end{keyidea}

Two features of this idea deserve emphasis, because they explain when \cbs
works well. First, the tree grows with the number of \emph{conflicts}, not
with the number of agents: ten drones that never meet cost ten single-agent
searches and one node. Second, the high level never looks inside a route. It
only asks the low level for the cheapest route under a list of forbidden
(cell, time) pairs, which is precisely the space-time \astar of
\cref{ch:ch04}. Any single-agent planner that returns optimal paths under
such constraints can serve as the low level, which is why \cbs generalises
easily to other graphs, action sets and cost functions.

% ---------------------------------------------------------------------
\section{Problem statement and notation}
\label{sec:ch09-problem}

We use the MAPF model of \cref{ch:ch07} (after Stern et al.~\cite{stern2019mapf}).
A graph $G = (V, E)$, in this chapter a 4-connected grid, is shared by $k$
agents $a_1, \dots, a_k$ with distinct starts $s_i$ and distinct goals
$\gamma_i$. Time is discrete. In one time step an agent moves along an edge
or waits in its cell. A \textbf{path} $\pi_i$ of agent $a_i$ is a sequence of
cells $\pi_i[0] = s_i, \pi_i[1], \dots, \pi_i[T_i] = \gamma_i$ with
consecutive cells equal or adjacent; the agent stays at $\gamma_i$ for all
$t \ge T_i$, so $\pi_i[t]$ is defined for every $t \ge 0$. Its \textbf{cost}
$\cost(\pi_i) = T_i$ is the time of the final arrival at the goal, and the
\textbf{sum of costs}\index{sum of costs} of a plan $\Pi = (\pi_1, \dots, \pi_k)$
is $\sumcost(\Pi) = \sum_{i=1}^{k} \cost(\pi_i)$. Two agents have a
\textbf{vertex conflict} $\langle a_i, a_j, v, t\rangle$ if
$\pi_i[t] = \pi_j[t] = v$, and an \textbf{edge conflict}
$\langle a_i, a_j, u, v, t\rangle$ if $\pi_i[t] = u$, $\pi_i[t+1] = v$,
$\pi_j[t] = v$ and $\pi_j[t+1] = u$.\index{conflict!vertex}\index{conflict!edge}
This is the swapping conflict of \cref{def:ch07-conflicts}; the \cbs
literature calls it an edge conflict because the constraint that forbids it
is an edge constraint. The same-direction edge conflict of \cref{ch:ch07}
needs no separate test here: by \cref{prop:ch07-hierarchy}(a) it implies a
vertex conflict.
Note that conflicts are checked for all $t$, including times after one of
the agents has arrived: a parked drone still occupies its goal cell. A
\textbf{solution} is a plan without conflicts; an \textbf{optimal} solution
minimises $\sumcost$, and $C^*$ denotes its cost.

\begin{definition}[Constraints]
\label{def:ch09-constraint}
A \textbf{vertex constraint}\index{constraint!vertex} $\langle a_i, v, t\rangle$
forbids agent $a_i$ to occupy cell $v$ at time $t$. An \textbf{edge
constraint}\index{constraint!edge} $\langle a_i, u, v, t\rangle$ forbids
$a_i$ to move from $u$ at time $t$ to $v$ at time $t+1$. A path $\pi_i$
\textbf{respects} a set $\mathcal{C}_i$ of constraints on $a_i$ if it
violates none of them, where a vertex constraint $\langle a_i, v, t\rangle$
with $v = \gamma_i$ and $t \ge \cost(\pi_i)$ counts as violated, because the
agent is parked at $\gamma_i$ at that time.
\end{definition}

\begin{definition}[Constraint-tree node]
\label{def:ch09-ctnode}
A node $N$ of the \textbf{constraint tree}\index{constraint tree!node} is a
triple $N = (N.\mathcal{C}, N.\pi, N.\cost)$ where $N.\mathcal{C} =
(N.\mathcal{C}_1, \dots, N.\mathcal{C}_k)$ assigns a set of constraints to
every agent, $N.\pi = (N.\pi_1, \dots, N.\pi_k)$ assigns to every agent a
path of \emph{minimum cost among the paths that respect}
$N.\mathcal{C}_i$, and $N.\cost = \sum_i \cost(N.\pi_i)$. The \textbf{root}
$R$ has $R.\mathcal{C}_i = \emptyset$ for every agent, so $R.\pi_i$ is an
unconstrained shortest path. A node is a \textbf{goal node} if $N.\pi$ has
no conflict. Every non-goal node has two \textbf{children}, obtained by
splitting one of its conflicts as follows:
\begin{equation}
\label{eq:ch09-split}
\begin{aligned}
\CbsSplit\bigl(\langle a_i, a_j, v, t\rangle\bigr) &= \set{\langle a_i, v, t\rangle,\ \langle a_j, v, t\rangle},\\
\CbsSplit\bigl(\langle a_i, a_j, u, v, t\rangle\bigr) &= \set{\langle a_i, u, v, t\rangle,\ \langle a_j, v, u, t\rangle}.
\end{aligned}
\end{equation}
The child that receives constraint $\kappa$ on agent $a_i$ has
$N'.\mathcal{C}_i = N.\mathcal{C}_i \cup \set{\kappa}$, the same constraint
sets as $N$ for all other agents, and the same paths as $N$ except
$N'.\pi_i$, which is recomputed.\index{conflict!splitting}
\end{definition}

The wording ``minimum cost among the paths that respect $N.\mathcal{C}_i$''
is the heart of the definition. It makes $N.\cost$ a quantity that the low
level can compute exactly, and it is what turns $N.\cost$ into a lower bound
on every solution that obeys $N$'s constraints (\cref{sec:ch09-properties}).
The problem the low level solves is therefore:

\begin{definition}[The low-level problem]
\label{def:ch09-lowlevel}
Given an agent $a_i$ and a finite set $\mathcal{C}_i$ of constraints on it,
find a path from $s_i$ to $\gamma_i$ of minimum cost that respects
$\mathcal{C}_i$, or report that none exists.
\end{definition}

This is the constrained space-time \astar of \cref{sec:ch04-spacetime} with the
constraint set in the role of the reservation table (\cref{ch:ch08}).
Three details of that
search matter enough to repeat them here. The state is a pair $(v, t)$; a
successor $(v', t+1)$ with $v' \in \set{v} \cup \Succ(v)$ is generated only
if neither $\langle a_i, v', t+1\rangle$ nor $\langle a_i, v, v', t\rangle$
is in $\mathcal{C}_i$. The heuristic is the static shortest-path distance
$d_i(v)$ to $\gamma_i$, computed once per agent by a backward breadth-first
search (or Dijkstra, \cref{ch:ch03}); it is consistent, so no state is
expanded twice. And the goal test is not ``$v = \gamma_i$'' but ``$v =
\gamma_i$ and $t > t_\gamma$''. This is the \textbf{goal-stay
test}\index{goal-stay test} of \cref{ch:ch08}, also called the
goal-occupied-later test\index{goal-occupied-later test}; here
\begin{equation}
\label{eq:ch09-tgamma}
t_\gamma = \max\set{t : \langle a_i, \gamma_i, t\rangle \in \mathcal{C}_i}
\quad (t_\gamma = -1 \text{ if no constraint mentions } \gamma_i),
\end{equation}
because a path that ends at time $t \le t_\gamma$ would park the agent on
its goal at the forbidden time. Since $\gcost(v,t) = t$ in space-time, the
search key is simply $\fcost(v,t) = t + d_i(v)$.

% ---------------------------------------------------------------------
\section{The algorithm}
\label{sec:ch09-algorithm}

\Cref{alg:ch09-cbs} is the high level and \cref{alg:ch09-lowlevel} the low
level. Together they are a complete description of \cbs as published by
Sharon et al.~\cite{sharon2015cbs}; the only freedom left is which conflict
to split when a node has several, and how to break ties in \Open.

\begin{algorithm}[htb]
\caption{Conflict-based search: the high level}
\label{alg:ch09-cbs}
\KwIn{graph $G$; agents $a_1, \dots, a_k$ with starts $s_i$ and goals $\gamma_i$}
\KwOut{a conflict-free plan $\Pi = (\pi_1, \dots, \pi_k)$ of minimum sum of costs, or \emph{failure}}
\Fn{\CbsHighLevel{$G$, $s_{1..k}$, $\gamma_{1..k}$}}{
  \ForEach{agent $a_i$}{$R.\mathcal{C}_i \gets \emptyset$; $R.\pi_i \gets$ \CbsLowLevel{$a_i$, $\emptyset$}\label{alg:ch09-cbs:root}\;}
  \lIf{some $R.\pi_i$ is failure}{\KwRet failure \tcp*[f]{an agent cannot reach its goal at all}}
  $R.\cost \gets \sum_i \cost(R.\pi_i)$; \Open $\gets \set{R}$\label{alg:ch09-cbs:open}\;
  \While{\Open $\neq \emptyset$}{
    $N \gets$ pop the node with the smallest $N.\cost$ from \Open (ties: fewer conflicts first)\label{alg:ch09-cbs:pop}\;
    $c \gets$ \CbsFirstConflict{$N.\pi$}\label{alg:ch09-cbs:validate}\;
    \lIf{$c =$ \KwNil}{\KwRet $N.\pi$ \tcp*[f]{goal node: conflict-free, hence optimal}\label{alg:ch09-cbs:goal}}
    \ForEach{constraint $\kappa \in$ \CbsSplit{$c$} \emph{(\cref{eq:ch09-split})}}{\label{alg:ch09-cbs:split}
      $a_i \gets$ the agent named in $\kappa$\;
      $N' \gets$ a copy of $N$; $N'.\mathcal{C}_i \gets N.\mathcal{C}_i \cup \set{\kappa}$\label{alg:ch09-cbs:child}\;
      $N'.\pi_i \gets$ \CbsLowLevel{$a_i$, $N'.\mathcal{C}_i$}\tcp*{replan only the constrained agent}\label{alg:ch09-cbs:replan}
      \If{$N'.\pi_i \neq$ failure}{
        $N'.\cost \gets \sum_j \cost(N'.\pi_j)$; insert $N'$ into \Open\label{alg:ch09-cbs:push}\;
      }
    }
  }
  \KwRet failure\;
}
\end{algorithm}

\begin{algorithm}[htb]
\caption{The low level: space-time \astar for one agent under its constraints}
\label{alg:ch09-lowlevel}
\KwIn{agent $a_i$ with start $s_i$ and goal $\gamma_i$; constraint set $\mathcal{C}_i$; static distances $d_i(v)$ to $\gamma_i$}
\KwOut{a minimum-cost path from $s_i$ to $\gamma_i$ that respects $\mathcal{C}_i$, or \emph{failure}}
\Fn{\CbsLowLevel{$a_i$, $\mathcal{C}_i$}}{
  $t_\gamma \gets$ largest $t$ with $\langle a_i, \gamma_i, t\rangle \in \mathcal{C}_i$, or $-1$\label{alg:ch09-lowlevel:tgamma}\;
  $T_{\max} \gets H + 1 + D$, with $H$ the largest time in $\mathcal{C}_i$ (or $-1$) and $D = \max_v d_i(v)$\label{alg:ch09-lowlevel:horizon}\;
  $\gcost(s_i, 0) \gets 0$; \Open $\gets \set{(s_i, 0)}$ with key $(d_i(s_i), 0)$; \Closed $\gets \emptyset$\;
  \While{\Open $\neq \emptyset$}{
    $(v, t) \gets$ pop the state with the smallest key $(\fcost, -\gcost)$ from \Open\label{alg:ch09-lowlevel:pop}\;
    \lIf{$(v, t) \in$ \Closed}{\KwContinue}
    add $(v, t)$ to \Closed\;
    \lIf{$v = \gamma_i$ \KwAnd $t > t_\gamma$}{\KwRet the path read off the parent pointers\label{alg:ch09-lowlevel:goal}}
    \lIf{$t \ge T_{\max}$}{\KwContinue\label{alg:ch09-lowlevel:cut}}
    \ForEach{$v' \in \set{v} \cup \Succ(v)$ with $(v', t+1) \notin$ \Closed}{\label{alg:ch09-lowlevel:succ}
      \lIf{$\langle a_i, v', t+1\rangle \in \mathcal{C}_i$ \KwOr $\langle a_i, v, v', t\rangle \in \mathcal{C}_i$}{\KwContinue\label{alg:ch09-lowlevel:forbid}}
      \If{$t + 1 < \gcost(v', t+1)$}{
        $\gcost(v', t+1) \gets t + 1$; parent$(v', t+1) \gets (v, t)$\;
        push $(v', t+1)$ into \Open with key $(t + 1 + d_i(v'),\, -(t+1))$\label{alg:ch09-lowlevel:push}\;
      }
    }
  }
  \KwRet failure\;
}
\end{algorithm}

\subsection{The high level, line by line}
\label{sec:ch09-highlevel}

Line~\ref{alg:ch09-cbs:root} builds the root: one unconstrained low-level
search per agent. If an agent cannot reach its goal even alone, no plan
exists and \cbs stops at once. Line~\ref{alg:ch09-cbs:open} puts the root
into \Open, a priority queue keyed on $N.\cost$.\index{CBS!high level}

Line~\ref{alg:ch09-cbs:pop} is the best-first choice: the node with the
smallest sum of costs is expanded. Among nodes of equal cost, expanding the
one with fewer conflicts first is the tie-breaking rule that Sharon et
al.\ recommend~\cite{sharon2015cbs}; a node without conflicts is a goal
node, and one with few conflicts is likely to be close to one. The
tie-breaking rule affects only the effort, not the answer.

Line~\ref{alg:ch09-cbs:validate} \emph{validates} the node: it scans the
$k$ paths, time step by time step, for the earliest vertex or edge conflict
(\cref{ch:ch07} gives the pairwise detection routine). The scan runs up to
the largest path length and treats every agent as parked at its goal after
its arrival. If no conflict exists, line~\ref{alg:ch09-cbs:goal} returns the
paths: the node is a goal node, and \cref{sec:ch09-properties} proves that
it is optimal. Otherwise the earliest conflict $c$ is split.

Lines~\ref{alg:ch09-cbs:split}--\ref{alg:ch09-cbs:push} generate the two
children. Each child copies the parent, so it inherits \emph{all} constraints
of the parent, and adds the single constraint $\kappa$ to the set of the
agent named in it (line~\ref{alg:ch09-cbs:child}). Only that agent is
replanned (line~\ref{alg:ch09-cbs:replan}): the other $k-1$ paths still
respect their unchanged constraint sets and are still of minimum cost, so
recomputing them would only waste time and could, with different
tie-breaking inside the low level, return different but equally good paths
and confuse a trace. If the low level fails, the child is discarded: no path
of $a_i$ satisfies its constraints, so no solution exists below that child.
Otherwise the child's cost is recomputed and the child joins \Open
(line~\ref{alg:ch09-cbs:push}).

Note what is \emph{not} in the loop. There is no closed list: the constraint
tree is a tree, and two nodes with the same constraint sets are possible but
rare (\cref{sec:ch09-implementation}). There is no goal test at generation
time: a child may be conflict-free but must wait in \Open until it is the
cheapest, because a cheaper node may still lead to a cheaper solution. And
there is no explicit stopping rule for unsolvable instances: if the instance
has no solution but every agent can reach its goal alone, the loop runs for
ever, generating ever more expensive nodes (\cref{sec:ch09-properties}).

\subsection{The low level, line by line}
\label{sec:ch09-lowlevel}

\Cref{alg:ch09-lowlevel} is \cref{sec:ch04-spacetime}'s space-time \astar written
against a constraint set.\index{CBS!low level} Line~\ref{alg:ch09-lowlevel:tgamma}
computes $t_\gamma$ of \cref{eq:ch09-tgamma}, the last time at which the goal
is forbidden. Line~\ref{alg:ch09-lowlevel:horizon} fixes the time horizon:
after the last constrained time $H$ the search is unconstrained, and from
any cell the goal is then at most $D$ steps away, so an optimal path never
needs more than $H + 1 + D$ steps. The horizon keeps the search finite when
no path exists, for instance when the constraints seal the agent in. (For
an edge constraint $\langle a_i, u, v, t\rangle$ the relevant time is
$t + 1$.)

Line~\ref{alg:ch09-lowlevel:pop} pops the state with the smallest
$\fcost = t + d_i(v)$, with ties broken towards the larger $t$ (the
``deeper'' state), exactly as in \cref{ch:ch04}. Line~\ref{alg:ch09-lowlevel:goal}
is the goal test with the goal-stay rule; a path that reaches
$\gamma_i$ at a time $t \le t_\gamma$ is not accepted, and the search
continues, typically by waiting somewhere and re-entering the goal later, or
by leaving the goal and coming back. Line~\ref{alg:ch09-lowlevel:succ}
generates the wait action and the moves; line~\ref{alg:ch09-lowlevel:forbid}
drops the successors that a vertex or edge constraint forbids. Everything
else is plain \astar with a consistent heuristic: because $d_i$ is
consistent, a state is expanded once, and the first path that passes the
goal test is a minimum-cost path among those that respect $\mathcal{C}_i$.

% ---------------------------------------------------------------------
\section{A worked example}
\label{sec:ch09-example}

\begin{example}[Two flight corridors that cross]
\label{ex:ch09-crossing}
\Cref{fig:ch09-example} (left) shows a $5 \times 3$ grid in which a
horizontal corridor (row~1) and a vertical corridor (column~2) cross in the
cell $(2,1)$; the four dark cells are buildings. Three drones fly:
$a_1$ from $(0,1)$ to $\gamma_1 = (4,1)$ along the horizontal corridor,
$a_2$ from $(2,0)$ up to $\gamma_2 = (2,2)$, and $a_3$ from $(2,2)$ down
through the crossing and then east to $\gamma_3 = (4,0)$. Cells are written
$(x, y)$ with $x$ the column and $y$ the row, $y$ growing upwards.
\end{example}

\begin{figure}[tb]
  \centering
  \inputfigure{ch09/example}
  \caption{The worked example. Left: the root of the constraint tree, in
    which every drone flies its own shortest path; $a_2$ and $a_3$ both want
    the crossing cell $(2,1)$ at $t = 1$ (hatched). Right: the optimal plan
    returned in node $N_3$, in which $a_2$ waits one step at its start and
    $a_1$ waits one step at $(1,1)$; the sum of costs rises from $10$ to
    $12$.}
  \label{fig:ch09-example}
\end{figure}

\paragraph{The root.} Every number below is produced by
\code{code/ch09\_cbs.py}, which prints this trace when run. Planned alone,
the drones take
\begin{align*}
R.\pi_1 &= (0,1),(1,1),(2,1),(3,1),(4,1), &
R.\pi_2 &= (2,0),(2,1),(2,2), \\
R.\pi_3 &= (2,2),(2,1),(3,1),(4,1),(4,0),
\end{align*}
with costs $4$, $2$ and $4$, so $R.\cost = 10$. Validation finds exactly one
conflict: $a_2$ and $a_3$ are both in $(2,1)$ at $t = 1$, the vertex
conflict $\langle a_2, a_3, (2,1), 1\rangle$. (Agent $a_1$ follows one step
behind $a_3$ from $t=2$ on --- at $t=2$ it enters the cell $(2,1)$ that
$a_3$ left at $t=1$ --- but following is not a conflict in our model; this
is why the one-step delay of $a_3$ in $N_2$ makes the two paths coincide.)
The root is split into $N_1$ with the constraint $\langle a_2, (2,1), 1\rangle$
and $N_2$ with $\langle a_3, (2,1), 1\rangle$.

\paragraph{Depth one.} In $N_1$ only $a_2$ is replanned. The cheapest path
that is not in $(2,1)$ at $t=1$ waits one step: $N_1.\pi_2 =
(2,0),(2,0),(2,1),(2,2)$, cost $3$, so $N_1.\cost = 11$. But now $a_2$ is in
the crossing at $t = 2$, exactly when $a_1$ arrives there: $N_1$ has the
single conflict $\langle a_1, a_2, (2,1), 2\rangle$. In $N_2$ only $a_3$ is
replanned; it waits one step at its start, $N_2.\pi_3 =
(2,2),(2,2),(2,1),(3,1),(4,1),(4,0)$, cost $5$, and $N_2.\cost = 11$ as
well. This node is worse than it looks: $a_2$ now moves up into $(2,2)$ at
$t=1$ while $a_3$ moves down into $(2,1)$, an edge conflict
$\langle a_2, a_3, (2,1)\to(2,2), 1\rangle$, and from $t=2$ on $a_3$ occupies
the same cells as $a_1$ at the same times, three more vertex conflicts.
\Open now holds $N_1$ and $N_2$, both of cost $11$; the tie-breaking rule
picks $N_1$ (one conflict) before $N_2$ (four).

\paragraph{Depth two.} $N_1$ is split on $\langle a_1, a_2, (2,1), 2\rangle$.
Child $N_3$ adds $\langle a_1, (2,1), 2\rangle$: $a_1$ waits once at
$(1,1)$, $N_3.\pi_1 = (0,1),(1,1),(1,1),(2,1),(3,1),(4,1)$, cost $5$, total
$12$, and the three paths are conflict-free. Child $N_4$ adds
$\langle a_2, (2,1), 2\rangle$ to $a_2$'s set, which now holds two
constraints; $a_2$ waits twice, cost $4$, total $12$, also conflict-free.
Both children go into \Open. The high level next pops $N_2$, because its
cost $11$ is still smaller than $12$: \emph{a conflict-free node is not
returned until it is the cheapest node in} \Open. $N_2$ is split on its
edge conflict into $N_5$ (constraint $\langle a_2, (2,1)\to(2,2), 1\rangle$,
cost $12$, five conflicts) and $N_6$ ($\langle a_3, (2,2)\to(2,1), 1\rangle$,
cost $12$, two conflicts). Now \Open holds $N_3, N_4, N_5, N_6$, all of cost
$12$; $N_3$ has no conflict and was generated first, so it is popped,
validated, and returned. \Cref{fig:ch09-tree} draws the tree and
\cref{tab:ch09-ct} lists every node; \cref{tab:ch09-timeline} shows the
returned plan.

\begin{figure}[tb]
  \centering
  \inputfigure{ch09/tree}
  \caption{The constraint tree of \cref{ex:ch09-crossing}. Every child adds
    exactly one constraint to its parent's sets ($+$) and replans one agent.
    Costs never decrease downwards. Nodes $N_4$, $N_5$ and $N_6$ were
    generated but never expanded; $N_4$ is a second optimal solution.}
  \label{fig:ch09-tree}
\end{figure}

\begin{table}[tb]
  \centering\small
  \caption{The constraint-tree nodes of \cref{ex:ch09-crossing} in the order
    of generation. A child inherits its parent's constraints and adds the one
    listed; $\pi_1$, $\pi_2$, $\pi_3$ are the costs of the three paths.
    ``Conflicts'' is the number of conflicts among the node's paths; the last
    column names the conflict that is split when the node is expanded.}
  \label{tab:ch09-ct}
  \footnotesize\setlength{\tabcolsep}{4pt}
  \begin{tabular}{@{}lllcccccl@{}}
  \toprule
  Node & Parent & Constraint added & $\pi_1$ & $\pi_2$ & $\pi_3$ & Total & Conflicts & Expansion \\
  \midrule
  $N_0$ & -- & none & 4 & 2 & 4 & 10 & 1 & 1st: split $\langle a_2,a_3,(2,1),1\rangle$\\
  $N_1$ & $N_0$ & $\langle a_2,(2,1),1\rangle$ & 4 & 3 & 4 & 11 & 1 & 2nd: split $\langle a_1,a_2,(2,1),2\rangle$\\
  $N_2$ & $N_0$ & $\langle a_3,(2,1),1\rangle$ & 4 & 2 & 5 & 11 & 4 & 3rd: split $\langle a_2,a_3,(2,1)\!\to\!(2,2),1\rangle$\\
  $N_3$ & $N_1$ & $\langle a_1,(2,1),2\rangle$ & 5 & 3 & 4 & 12 & 0 & 4th: goal node, returned\\
  $N_4$ & $N_1$ & $\langle a_2,(2,1),2\rangle$ & 4 & 4 & 4 & 12 & 0 & left in \Open (also optimal)\\
  $N_5$ & $N_2$ & $\langle a_2,(2,1)\!\to\!(2,2),1\rangle$ & 4 & 3 & 5 & 12 & 5 & left in \Open\\
  $N_6$ & $N_2$ & $\langle a_3,(2,2)\!\to\!(2,1),1\rangle$ & 4 & 2 & 6 & 12 & 2 & left in \Open\\
  \bottomrule
  \end{tabular}
\end{table}

\begin{table}[tb]
  \centering\small
  \caption{The plan returned in $N_3$: the cell of every drone at every
    time step. $a_3$ passes the crossing first, $a_2$ follows, and $a_1$
    crosses last; after its arrival a drone stays at its goal.}
  \label{tab:ch09-timeline}
  \begin{tabular}{@{}lcccccc@{}}
  \toprule
  $t$ & 0 & 1 & 2 & 3 & 4 & 5\\
  \midrule
  $a_1$ & $(0,1)$ & $(1,1)$ & $(1,1)$ wait & $(2,1)$ & $(3,1)$ & $(4,1)$ goal\\
  $a_2$ & $(2,0)$ & $(2,0)$ wait & $(2,1)$ & $(2,2)$ goal & $(2,2)$ & $(2,2)$\\
  $a_3$ & $(2,2)$ & $(2,1)$ & $(3,1)$ & $(4,1)$ & $(4,0)$ goal & $(4,0)$\\
  \bottomrule
  \end{tabular}
\end{table}

\paragraph{What the example shows.} \cbs expanded $4$ nodes, generated
$7$, and called the low level $9$ times (three times for the root, once per
child); the optimal cost is $12$, which the brute-force search over the joint
state space in \code{ch09\_cbs.py} confirms. Four lessons are worth
recording. (i)~Resolving one conflict can create another: the wait of $a_2$
in $N_1$ produced the conflict with $a_1$, a drone that was not involved at
the root. (ii)~Constraints accumulate: $N_4$ holds two constraints on $a_2$,
and its path respects both. (iii)~The whole subtree below $N_2$, in which
$a_3$ yields, costs at least $12$; \cbs generated its two children but never
had to expand them, because $N_3$ was found at the same cost first, and it
never had to look deeper, because costs cannot decrease. (iv)~Two optimal
solutions exist ($N_3$ and $N_4$); which one is returned depends only on
tie-breaking. Finally, every conflict split here was \emph{cardinal}: both
children were more expensive than their parent. \Cref{sec:ch09-variants}
explains why this is the case one wants, and why the variant that splits
cardinal conflicts first produces the same tree on this instance.

% ---------------------------------------------------------------------
\section{Properties: correctness, optimality, complexity}
\label{sec:ch09-properties}

Throughout this section the graph is finite, every agent can reach its
goal alone, and the low level is \cref{alg:ch09-lowlevel}, which returns a
minimum-cost path respecting its constraints whenever one exists
(\cref{ch:ch04}, together with the horizon argument of
\cref{sec:ch09-lowlevel}). Two lemmas carry the entire theory.

\begin{lemma}[A split loses no solution]
\label{thm:ch09-split}
Let $N$ be a constraint-tree node whose paths have a conflict $c$ between
$a_i$ and $a_j$, and let $N_i'$ and $N_j'$ be the two children obtained by
splitting $c$. Every solution $\Pi$ whose paths respect the constraints of
$N$ respects the constraints of $N_i'$ or those of $N_j'$ (or both).
\end{lemma}
\begin{proof}
Suppose $c = \langle a_i, a_j, v, t\rangle$ is a vertex conflict. If
$\pi_i[t] = v$ and $\pi_j[t] = v$, then $\Pi$ has a vertex conflict at
$(v,t)$ and is not a solution. Hence $\pi_i[t] \neq v$ or $\pi_j[t] \neq v$.
In the first case $\pi_i$ respects $\langle a_i, v, t\rangle$, the only
constraint of $N_i'$ that is not already a constraint of $N$, so $\Pi$
respects all constraints of $N_i'$; in the second case the same holds for
$N_j'$. (The convention that a parked agent occupies its goal is used here:
if $\pi_i$ has arrived at $\gamma_i = v$ before $t$, then $\pi_i[t] = v$ and
$\Pi$ violates the constraint, exactly as it would conflict in reality.) If
$c = \langle a_i, a_j, u, v, t\rangle$ is an edge conflict, the children
forbid $a_i$ to move $u \to v$ and $a_j$ to move $v \to u$ at time $t$. A
solution cannot make both moves, because that is an edge conflict, so it
respects at least one of the two constraints. \qedhere
\end{proof}

\begin{lemma}[The cost of a node is a lower bound]
\label{thm:ch09-lowerbound}
Let $N$ be a constraint-tree node and $\Pi$ any plan whose paths respect the
constraints of $N$. Then $\sumcost(\Pi) \ge N.\cost$. In particular, every
descendant $N'$ of $N$ satisfies $N'.\cost \ge N.\cost$: costs never decrease
along a branch of the tree.
\end{lemma}
\begin{proof}
By \cref{def:ch09-ctnode}, $N.\pi_i$ has minimum cost among the paths of
$a_i$ that respect $N.\mathcal{C}_i$, and $\pi_i$ is such a path, so
$\cost(\pi_i) \ge \cost(N.\pi_i)$ for every $i$; summing over $i$ gives the
first claim. For the second, the constraint sets of a descendant contain
those of $N$, so $N'.\pi$ respects the constraints of $N$, and the first
claim applies with $\Pi = N'.\pi$. \qedhere
\end{proof}

\begin{lemma}[Invariant of the high level]
\label{thm:ch09-invariant}
Suppose the instance is solvable and let $\Pi^*$ be an optimal solution,
$C^* = \sumcost(\Pi^*)$. At the start of every iteration of the loop of
\cref{alg:ch09-cbs}, \Open contains a node $N$ whose constraints are all
respected by $\Pi^*$, and every such node has $N.\cost \le C^*$.
\end{lemma}
\begin{proof}
The bound $N.\cost \le C^*$ is \cref{thm:ch09-lowerbound} applied to
$\Pi = \Pi^*$. For the existence, argue by induction over the iterations.
Initially \Open holds the root, which has no constraints. Consider an
iteration that pops a node $N$. If $N$ is not one of the nodes respected by
$\Pi^*$, those nodes stay in \Open. If it is, and it is a goal node, the
algorithm returns and there is no next iteration. Otherwise $N$ has a
conflict $c$, and by \cref{thm:ch09-split} $\Pi^*$ respects the constraints
of at least one child $N'$. The path $\pi_i^*$ of the agent constrained in
$N'$ respects $N'.\mathcal{C}_i$, so a path exists and the low level does
not fail; $N'$ is generated and inserted into \Open. \qedhere
\end{proof}

\begin{theorem}[Optimality of \cbs]
\label{thm:ch09-optimal}
If \cref{alg:ch09-cbs} returns a plan $\Pi$, then $\Pi$ is a solution and
$\sumcost(\Pi) = C^*$.\index{CBS!optimality}
\end{theorem}
\begin{proof}
The plan is returned at line~\ref{alg:ch09-cbs:goal} only after
line~\ref{alg:ch09-cbs:validate} found no conflict, so $\Pi$ is a solution
and $\sumcost(\Pi) \ge C^*$. Let $N$ be the node that was popped; $N.\cost =
\sumcost(\Pi)$. At that moment \Open contained, by
\cref{thm:ch09-invariant}, a node of cost at most $C^*$, and $N$ was a node
of minimum cost in \Open, so $N.\cost \le C^*$. Hence
$\sumcost(\Pi) = C^*$. \qedhere
\end{proof}

\begin{theorem}[Completeness of \cbs on solvable instances]
\label{thm:ch09-complete}
If the instance has a solution, \cref{alg:ch09-cbs} terminates and returns
one.\index{CBS!completeness}
\end{theorem}
\begin{proof}
Let $C^*$ be the optimal cost. We first show that only finitely many
constraint-tree nodes have cost at most $C^*$. Let $N$ be such a node. By
\cref{thm:ch09-lowerbound} all its ancestors have cost at most $C^*$ too.
Every constraint on the branch from the root to $N$ was created by splitting
a conflict of some ancestor $M$ at a time $t$; at that time at least one of
the two agents had not yet arrived (two parked agents cannot share a cell,
because goals are distinct, and an edge conflict involves two moving
agents), so $t \le \max_i \cost(M.\pi_i) \le M.\cost \le C^*$. Moreover the
constraint was not already in $M$'s sets, because $M$'s path violated it
while respecting $M$'s constraints. The constraints along the branch are
therefore distinct elements of the finite set of constraints with time at
most $C^* + 1$, whose size is at most $\Delta = k\,(\abs{V} + 2\abs{E})\,(C^* + 2)$.
So the depth of $N$ is at most $\Delta$, and since every node has at most
two children, there are fewer than $2^{\Delta + 1}$ nodes of cost at most
$C^*$.

Now run the algorithm. By \cref{thm:ch09-invariant}, \Open is never empty,
so the loop can only end by returning at line~\ref{alg:ch09-cbs:goal}.
Every popped node has minimum cost in \Open and \Open always contains a
node of cost at most $C^*$, so every popped node has cost at most $C^*$.
Each node is popped at most once, and there are finitely many nodes of cost
at most $C^*$, so the loop performs finitely many iterations and must
return. By \cref{thm:ch09-optimal} the returned plan is an optimal solution.
\qedhere
\end{proof}

\begin{remark}[Unsolvable instances]
\label{rem:ch09-unsolvable}
If every agent can reach its goal alone but no conflict-free plan exists
(two agents in a corridor that must swap, for instance), the invariant
fails and \cbs keeps generating ever more expensive nodes for ever. Sharon
et al.~\cite{sharon2015cbs} state completeness for solvable instances only.
In practice one imposes a limit on time or on the number of expanded nodes,
or first runs a complete polynomial-time solver such as Push-and-Rotate
(\cref{ch:ch11}) to decide solvability and to obtain an upper bound on
$C^*$, beyond which the constraint tree need not be explored.
\end{remark}

\subsection{What the running time depends on}
\label{sec:ch09-complexity}

One iteration of the high level costs one validation, $\bigO{k^2 T}$ for
paths of length up to $T$ when done pairwise (or $\bigO{kT}$ with a hash
table of occupied cells per time step), plus two low-level searches, each
$\bigO{\abs{V}\,T_{\max} \log(\abs{V}\,T_{\max})}$ for the space-time
\astar. The total time is this per-node cost multiplied by the number of
expanded constraint-tree nodes, and that number is the crux. The proof of
\cref{thm:ch09-complete} bounds it by $2^{\Delta + 1}$, where $\Delta$ counts
the constraints that could appear; in words, \emph{the effort of \cbs is
exponential in the number of conflicts it has to resolve, and only
polynomial in the number of agents for a fixed number of conflicts}.
Sharon et al.~\cite{sharon2015cbs} phrase the same conclusion as a
comparison with \astar on the joint state space (with the operator
decomposition and independence detection of Standley~\cite{standley2010finding},
see \cref{ch:ch11}): joint-space search is exponential in $k$ regardless of
how the agents interact, \cbs is exponential in the depth of the constraint
tree regardless of $k$.\index{CBS!complexity}

\paragraph{Corridors against open rooms.} Which of the two is better on a
given instance depends on how many alternative paths a conflict leaves open.
Consider first a \emph{bottleneck}: two agents that must pass, one after the
other, through a gap of one cell (\cref{exr:ch09-handtree}). Each agent has
a single shortest path, one constraint delays one agent by one step, and
the constraint tree has three nodes; joint-space \astar, in contrast, has to
explore the product of the two agents' approaches to the gap. \cbs wins
easily, and the more free space surrounds the bottleneck, the larger its
advantage. Now consider an \emph{open room}: two agents cross an obstacle-free
rectangle in perpendicular directions. Both have many shortest paths of
equal cost, and, as \cref{sec:ch09-symmetry} and \cref{exr:ch09-rectangle}
show, every pair of them conflicts. A constraint on one shortest path
merely diverts the agent to another shortest path that conflicts again at a
different cell, at the same cost, so the constraint tree fills a whole level
of equal-cost nodes before the cost can rise by one. On the $5 \times 5$
room of \cref{fig:ch09-symmetry} the self-test of \code{ch09\_cbs.py}
records $15$ expanded nodes for a single cost increase with only two agents;
the number grows exponentially with the size of the rectangle, whereas
joint-space \astar simply steps around the encounter. The worst case for
\cbs is a small, crowded map in which agents must shuffle past each other.
The self-test contains a $4 \times 3$ maze with $8$ free cells and three
agents whose optimal sum of costs is $24$: Dijkstra on the joint state space
finds it in $0.01$~s, while \cbs was still without a conflict-free node
after the $2\,000$ constraint-tree nodes ($0.3$~s) that the self-test
allows it, because every extra
unit of cost can be spent as a wait of one step by one of several agents
at one of several times, and each of these choices is a separate branch.
\Cref{tab:ch09-compare} summarises the comparison; the improvements of
\cref{sec:ch09-variants} attack precisely the open-room and corridor cases.

\begin{table}[tb]
  \centering\small
  \caption{Exact and heuristic MAPF solvers seen so far. ``Effort grows
    with'' names the quantity that drives the exponential part of the
    running time.}
  \label{tab:ch09-compare}
  \footnotesize
  \begin{tabularx}{\textwidth}{@{}lccYY@{}}
  \toprule
  Method & Complete & Optimal & Effort grows with & Best on\\
  \midrule
  Joint-space \astar (\cref{ch:ch07,ch:ch11}) & \checkmark & \checkmark & number of agents $k$ & few, tightly coupled agents\\
  Prioritized planning (\cref{ch:ch08}) & no & no & nothing exponential & many agents, sparse interaction, no guarantees needed\\
  \cbs (this chapter) & solvable instances & \checkmark & depth of the constraint tree (conflicts to resolve) & many agents with few, local conflicts\\
  \ecbs (\cref{ch:ch10}) & solvable instances & $w$-bounded & the same, with far fewer nodes & the same, at larger scale\\
  \bottomrule
  \end{tabularx}
\end{table}

% ---------------------------------------------------------------------
\section{Variants and extensions}
\label{sec:ch09-variants}

Plain \cbs, as defined so far, is a clean algorithm but a slow solver: on
open maps it spends most of its time generating equal-cost nodes that
resolve nothing. The improvements below, all of which keep optimality,
attack this waste at four points: \emph{which} conflict to split, whether
to split at all, how to \emph{guide} the high level, and how to recognise
whole families of conflicts at once. The bounded-suboptimal relaxation of
\cbs, \ecbs, is the subject of \cref{ch:ch10}.

\subsection{Cardinal, semi-cardinal and non-cardinal conflicts}
\label{sec:ch09-cardinal}

\begin{definition}[Conflict types]
\label{def:ch09-cardinal}
Let $c$ be a conflict of node $N$ and let $N_i'$ and $N_j'$ be the children
that splitting $c$ would produce. The conflict is
\textbf{cardinal}\index{conflict!cardinal} if $N_i'.\cost > N.\cost$ and
$N_j'.\cost > N.\cost$, \textbf{semi-cardinal}\index{conflict!semi-cardinal}
if exactly one child is more expensive than $N$, and
\textbf{non-cardinal}\index{conflict!non-cardinal} if neither is.
\end{definition}

Boyarski et al.~\cite{boyarski2015icbs} observed that the order in which
conflicts are split changes the size of the tree dramatically, and that
splitting a non-cardinal conflict is almost always wasted work: both
children keep the parent's cost, so \Open receives two nodes that are as
good as their parent and still have conflicts, and the level of equal-cost
nodes doubles. Splitting a cardinal conflict, in contrast, raises the cost in
both children, so the search makes progress towards $C^*$ with every split.
Their algorithm, improved \cbs (ICBS)\index{ICBS}\index{CBS!improved},
therefore splits a cardinal conflict if the node has
one, else a semi-cardinal one, else any conflict. \Cref{tab:ch09-conflict-types}
summarises the three cases; in \cref{ex:ch09-crossing} every split was
cardinal, which is why ICBS builds the same tree there.

\begin{table}[tb]
  \centering\small
  \caption{The three conflict types of ICBS and what a split achieves.}
  \label{tab:ch09-conflict-types}
  \begin{tabular}{@{}llll@{}}
  \toprule
  Type & Children's costs & MDD test at level $t$ & Split it? \\
  \midrule
  cardinal & both $> N.\cost$ & single cell in both MDDs & yes, first: the cost rises\\
  semi-cardinal & exactly one $> N.\cost$ & single cell in one MDD & second choice\\
  non-cardinal & both $= N.\cost$ & alternatives in both MDDs & avoid; try a bypass\\
  \bottomrule
  \end{tabular}
\end{table}

Recognising the type without generating the children uses the
\textbf{multi-valued decision diagram}\index{MDD} (MDD) of an agent,
introduced by Sharon et al.\ for the increasing cost tree
search~\cite{sharon2013icts}: for agent $a_i$ with current cost $c =
\cost(N.\pi_i)$, the MDD is the layered graph whose level $t$ contains every
cell that lies at time $t$ on \emph{some} path of cost $c$ from $s_i$ to
$\gamma_i$ respecting $N.\mathcal{C}_i$. Building it costs one breadth-first
search forwards and one backwards in space-time. A vertex conflict at
$(v, t)$ is cardinal exactly when $v$ is the only cell at level $t$ of both
MDDs~\cite{boyarski2015icbs}, semi-cardinal when it is the only cell in one of them, and
non-cardinal otherwise; edge conflicts are classified in the same way on
the MDD edges. \Cref{fig:ch09-mdd} shows both situations: the corridor
agents of the worked example have chain-shaped MDDs, so their conflict is
cardinal, whereas two agents crossing an open room have wide MDDs, and all
their conflicts are non-cardinal, the reason for the exponential trees of
\cref{sec:ch09-complexity}.

\begin{figure}[tb]
  \centering
  \inputfigure{ch09/mdd}
  \caption{Multi-valued decision diagrams. Left: at the root of
    \cref{ex:ch09-crossing} the MDDs of $a_2$ and $a_3$ are chains that
    share the cell $(2,1)$ at level $1$, so the conflict there is cardinal.
    Right: in an open $3 \times 3$ room the MDDs of two crossing agents are
    wide; the shared cells (red) have alternatives on both sides, so every
    conflict is non-cardinal and a split does not raise the cost.}
  \label{fig:ch09-mdd}
\end{figure}

\subsection{Bypassing a conflict}
\label{sec:ch09-bypass}

The second ICBS idea handles the non-cardinal conflicts that remain.
Suppose node $N$ is expanded and one of its children, say $N'$, has the same
cost as $N$ but fewer conflicts. Instead of inserting both children, ICBS
\emph{adopts} the paths of $N'$ into $N$ ($N.\pi \gets N'.\pi$) and puts $N$
back into \Open; no split takes place. This is the
\textbf{bypass}\index{bypass}\index{CBS!bypass}: the agent found an
alternative path of the same cost that avoids the conflict, so the tree need
not record the choice. The modified $N$ is still a valid constraint-tree
node in the sense of \cref{def:ch09-ctnode}, because $N'.\pi_i$ respects
$N.\mathcal{C}_i$ (it respects the larger set $N'.\mathcal{C}_i$) and has the
minimum cost under it, because $\cost(N'.\pi_i) = \cost(N.\pi_i)$, which
is the minimum under $N.\mathcal{C}_i$; hence
\cref{thm:ch09-split,thm:ch09-lowerbound} and
the optimality proof are untouched. If no child qualifies, $N$ is split as
usual. Bypasses are cheap, because the children had to be generated anyway,
and on open maps they remove most non-cardinal conflicts without growing
the tree.

\subsection{Meta-agent \cbs}
\label{sec:ch09-macbs}

When two agents keep conflicting, a long chain of one-cell constraints is
a poor way to coordinate them. Meta-agent \cbs (MA-CBS)\index{meta-agent
CBS}\index{CBS!meta-agent} by Sharon et al.~\cite{sharon2015cbs} counts, for
every pair of agents, the conflicts split between them along the current
branch; when the count exceeds a bound $B$, the two agents are
\emph{merged} into a meta-agent that the low level plans jointly, in the
joint space of the pair, with the operator-decomposition \astar of
Standley~\cite{standley2010finding}. The node is then replanned with the
meta-agent under the union of the two agents' external constraints, and the
tree continues to branch on conflicts between (meta-)agents. With $B = 0$
every pair is merged at its first conflict, so MA-CBS degenerates to
Standley's independence detection~\cite{standley2010finding}: each group of
interacting agents is solved by a joint-space \astar, while agents that
never conflict are still planned alone. With $B = \infty$ it is plain
\cbs. The bound therefore interpolates between the two
regimes of \cref{sec:ch09-complexity}: small $B$ for tightly coupled
groups in crowded rooms, large $B$ for loosely coupled swarms in corridors.

\subsection{Admissible heuristics for the high level}
\label{sec:ch09-heuristics}

The high level of \cref{alg:ch09-cbs} orders nodes by $N.\cost$ alone,
which is $\gcost$ without an $\hcost$: it is uniform-cost search over the
constraint tree. Felner et al.~\cite{felner2018adding} showed how to add an
admissible heuristic and run the high level as \astar (the family is called
CBSH, for \cbs with heuristics)\index{CBSH}\index{CBS!high-level heuristic}. The quantity to estimate
is how much more cost any solution below $N$ must pay, and cardinal
conflicts provide a lower bound. Build the \textbf{conflict graph} of $N$:
one vertex per agent and one edge per cardinal conflict. Every solution
below $N$ resolves every cardinal conflict, and resolving one forces at
least one of its two agents to pay at least one extra unit (that is what
cardinal means, in unit-cost grids). Set
\begin{equation}
\label{eq:ch09-hcg}
h_{\mathrm{CG}}(N) = \text{size of a minimum vertex cover of the conflict graph of } N
\end{equation}
and collect the agents whose cost is larger in a solution below $N$ than in
$N$. Because each cardinal conflict forces at least one of its two agents
into that set, the set is a vertex cover of the conflict graph, so it
contains at least $h_{\mathrm{CG}}(N)$ agents, each paying at least one
extra unit. Hence every solution below $N$ costs at least
$N.\cost + h_{\mathrm{CG}}(N)$: the heuristic $h_{\mathrm{CG}}$ is
admissible, and the high level can expand nodes in order of
$N.\cost + h_{\mathrm{CG}}(N)$. Minimum vertex cover is NP-hard in general,
but conflict graphs have few edges and a small branch-and-bound solves them
in microseconds. At the root of \cref{ex:ch09-crossing} the conflict graph
has the single edge $a_2 a_3$, so $h_{\mathrm{CG}}(N_0) = 1$ and the root's
key is $11$; in $N_1$ the conflict $\langle a_1, a_2, (2,1), 2\rangle$ is
cardinal, so $N_1$ has key $12$, the same as $N_2$ (cost $11$ plus one
cardinal edge) and the same as $N_3$ (cost $12$, no conflict). A tie-break
towards fewer conflicts then pops $N_3$ before $N_2$, and CBSH expands three
nodes instead of four. Later work strengthened the heuristic with
dependency graphs (an edge whenever the pair's MDDs admit no conflict-free
joint path) and weighted dependency graphs (the edge weight is the pair's
true extra cost, and a minimum weighted vertex cover is
taken)~\cite{li2019improved}; these are the heuristics of today's
reference implementations.

\subsection{Symmetry breaking: rectangles and corridors}
\label{sec:ch09-symmetry}

The exponential trees of \cref{sec:ch09-complexity} come from
\emph{symmetries}: families of shortest paths that are interchangeable for
the cost but conflict pairwise, so that one-cell constraints exclude them
one at a time. Li et al.~\cite{li2019symmetry} identified the most common
one on grids, the \textbf{rectangle symmetry}\index{symmetry!rectangle}
(\cref{fig:ch09-symmetry}, left): two agents whose shortest paths cross an
obstacle-free rectangle in perpendicular directions, each moving
monotonically towards its goal. \Cref{exr:ch09-rectangle} asks you to prove
that every shortest path of one agent conflicts with every shortest path of
the other. \textbf{Rectangle reasoning} detects such a pair from the MDDs
and the start--goal bounding boxes, and replaces the two one-cell
constraints of the split by two \textbf{barrier constraints}: the first
child forbids $a_1$ from every cell of the rectangle's exit border at the
time at which a shortest path would reach it, the second child does the
same for $a_2$. The barriers are a valid split in the sense of
\cref{thm:ch09-split}, since no solution can send both agents across the
rectangle at shortest-path speed, so optimality is preserved, and a whole
exponential family of conflicts is resolved by one split.

\begin{figure}[tb]
  \centering
  \inputfigure{ch09/symmetry}
  \caption{Two symmetries. Left: rectangle symmetry in an open $5 \times 5$
    room; one of the $5$ shortest paths of $a_1$ and one of the $10$ of
    $a_2$ are drawn, and every pair meets in the shaded rectangle. Right:
    corridor symmetry; $a_1$ and $a_2$ enter a corridor of length $5$ from
    opposite ends, one must wait outside until the other has come through,
    and plain \cbs turns each possible waiting time into a branch.}
  \label{fig:ch09-symmetry}
\end{figure}

The mirror image of the open room is the \textbf{corridor
symmetry}\index{symmetry!corridor} (\cref{fig:ch09-symmetry}, right): two
agents traverse a corridor of degree-two cells in opposite directions. One
of them must wait outside until the other has left, so the optimal cost
exceeds the root cost by roughly the length of the corridor, and every
intermediate cost level is filled with equal-cost nodes that differ only in
where and how long the agents wait. \textbf{Corridor reasoning}, introduced
with several related techniques by Li et al.~\cite{li2020pairwise}, computes
for each agent the earliest time at which it can leave the corridor if the
other agent goes first, and splits with two \emph{range constraints}: the
first child forbids $a_1$ from the far end of the corridor at all times up
to that bound, the second child forbids $a_2$ symmetrically. The same paper
treats the \emph{target symmetry}, in which an agent parked at its goal
blocks a corridor that another agent must use later, the situation of
\cref{exr:ch09-goalstay}. All these techniques share one principle, which
is also the principle behind barrier constraints and MA-CBS: \emph{any pair
of constraint sets such that every solution satisfies at least one of them
is a legal split}. \Cref{thm:ch09-split} is the only property a split has to
satisfy.

\subsection{An experiment on random grids}
\label{sec:ch09-experiment}

\Cref{fig:ch09-scaling} shows how \cbs scales on random $8 \times 8$ grids
with $15\,\%$ obstacles as the number of agents grows from $2$ to $12$. For
each $k$ the script \code{code/figures/gen\_ch09\_scaling.py} draws $12$
instances with random distinct starts and goals and runs two solvers from
\code{ch09\_cbs.py} on each, with a limit of $5$~s and $20\,000$ expanded
nodes per run: plain \cbs (split the earliest conflict, ties on fewer
conflicts) and \cbs with cardinal conflicts first and the bypass, labelled
ICBS. The success rate is the fraction of instances solved within the
limits; the node counts are averaged over the solved instances only.

\begin{figure}[tb]
  \centering
  \inputfigure{ch09/scaling}
  \caption{\cbs on random $8 \times 8$ grids ($15\,\%$ obstacles, $12$
    instances per $k$, $5$~s and $20\,000$ CT nodes per run). Left: fraction
    of instances solved. Right: mean and maximum number of expanded
    constraint-tree nodes on the solved instances, logarithmic scale. Both
    curves grow exponentially with $k$ on this small map, because the number
    of conflicts does; cardinal-first splitting and bypasses divide the node
    count by three to seven and raise the success rate.}
  \label{fig:ch09-scaling}
\end{figure}

Up to six agents every instance is solved in a few milliseconds, with
$1.3$, $4.0$ and $7.8$ expanded nodes on average for $k = 2, 4, 6$. From
eight agents on the tree explodes on some instances: plain \cbs solves
$92\,\%$ of the instances with $k = 8$, $75\,\%$ with $k = 10$ and $58\,\%$
with $k = 12$, and the mean number of expanded nodes on the solved instances
grows to $55$, $241$ and $622$, with maxima of $441$, $1\,061$ and $2\,398$.
The gap between mean and maximum is the signature of \cref{sec:ch09-complexity}:
most instances are loosely coupled and cheap, a few contain a crowded room
or a corridor and dominate the running time. Cardinal-first splitting with
bypasses cuts the mean node count to $15$, $49$ and $92$ for $k = 8, 10,
12$ and the maxima to $66$, $199$ and $308$; it solves all instances with
$k = 8$ and $83\,\%$ and $75\,\%$ for $k = 10$ and $12$. Its price is
visible in the low-level statistics: to classify a conflict, our
implementation generates the two children of every candidate conflict in
time order until it finds a cardinal one, so it performs many more
low-level expansions per constraint-tree node: at $k = 12$ a solved run
costs $23\,600$ low-level state expansions against $17\,700$ for plain
\cbs, although it expands only $92$ instead of $622$ CT nodes; an MDD-based
classifier avoids this. Twelve instances per
setting make the success rates coarse (one instance is $8\,\%$), and the
absolute times are those of pure \python, but the shape of the curves is
the one reported by Sharon et al.\ and Boyarski et al.\ on benchmark maps:
the number of agents that \cbs handles is set by the number of conflicts on
the map, and the ICBS improvements move the wall by a few agents. Larger
teams need the bounded-suboptimal \ecbs of \cref{ch:ch10}.

% ---------------------------------------------------------------------
\section{Implementation notes}
\label{sec:ch09-implementation}

\code{code/ch09\_cbs.py} implements everything in this chapter in about
$600$ lines, of which some $110$ are the self-tests: the grid and the backward breadth-first distances, the low
level, conflict detection, the high level with the tie-breaking on
conflicts, the cardinal-first and bypass options, a brute-force optimal
solver on the joint state space for certification, and the worked example.
\Cref{lst:ch09-lowlevel} is the low level and \cref{lst:ch09-highlevel} the
body of the high level. The remarks below are the decisions that matter.

\paragraph{Constraint sets.} A node stores one immutable set of vertex
constraints $(v, t)$ and one of edge constraints $(u, v, t)$ per agent.
Immutability (\code{frozenset}) makes ``a child inherits all constraints of
its parent and adds one'' a single union, and it makes aliasing bugs
impossible: a child can never modify its parent's sets. For very deep trees
one stores only the added constraint and a parent pointer, and rebuilds the
set on demand, or uses a persistent set structure; the same applies to the
$k$ paths, of which only one differs from the parent.

\paragraph{The heuristic of the low level.} The static distances $d_i(v)$
to $\gamma_i$ are computed once per agent by a backward breadth-first
search and shared by every low-level call for that agent; the distances
also give $D$ for the horizon and tell at once whether an agent can reach
its goal at all. With a consistent heuristic no state is expanded twice, so
the closed set is a plain \code{set} of $(v, t)$ pairs.

\paragraph{Which conflict to split.} The default is the earliest conflict
in time, the choice of Sharon et al.: resolving an early conflict changes
the paths from that time on and often removes later conflicts as a side
effect, whereas resolving a late conflict first is frequently undone by the
resolution of an earlier one. The option \code{split="cardinal"} implements
ICBS's preference by generating the children of the conflicts in time
order and stopping at the first cardinal one; the generated children are
reused when the node is expanded. This is exact but costs extra low-level
searches (\cref{sec:ch09-experiment}); a production solver classifies
conflicts with MDDs instead.

\paragraph{Tie-breaking.} The heap key is
$(N.\cost,\ \text{number of conflicts},\ \text{node id})$. Fewer conflicts
first is the rule of \cref{alg:ch09-cbs}; the id makes the order
deterministic (first generated, first expanded), which the worked example
and the self-test rely on. Preferring \emph{deeper} nodes among equal keys
is another common choice: it dives towards a goal node instead of sweeping
a level.

\paragraph{Duplicate nodes.} Two branches can reach the same constraint
sets, for instance by splitting two independent conflicts in the two
possible orders. Duplicate detection by hashing the constraint sets is easy
to add, but such coincidences are uncommon; the implementation here does not
detect duplicates, and the
proofs of \cref{sec:ch09-properties} do not need it.

\paragraph{Validation and limits.} Conflict detection walks the time axis
once, comparing every pair of agents, and treats an agent as parked at its
goal after its arrival; on large teams a dictionary from cells to agents per
time step replaces the pairwise loop. The solver takes a time limit and a
node limit and returns \code{None} when it hits one, which is the only way
to stop on an unsolvable instance (\cref{rem:ch09-unsolvable}). The
function \code{validate} re-checks a returned plan independently (start,
goal, adjacency of consecutive cells, obstacles, conflicts); running it on
every answer during development catches most bugs of the kinds described in
the pitfalls below.

\begin{lstlisting}[caption={The low level: space-time \astar under constraints, with the goal-stay test and the horizon $T_{\max} = H + 1 + D$ (from \code{code/ch09\_cbs.py}).},label={lst:ch09-lowlevel}]
def low_level(grid: Grid, start: Cell, goal: Cell, cons: Constraints,
              dist: Dict[Cell, int], stats: Optional[Stats] = None) -> Optional[Path]:
    """Space-time A* for one agent: a shortest path from ``start`` to
    ``goal`` that respects ``cons``; ``dist`` are true distances to the goal.
    The agent may finish at (goal, t) only if no vertex constraint touches
    the goal at a time >= t, because it stays at the goal for ever after."""
    if stats is not None:
        stats.low_level_calls += 1
    if start not in dist:
        return None
    goal_block = max((t for v, t in cons.vertex if v == goal), default=-1)
    horizon = cons.last_time() + 1 + max(dist.values())     # T_max = H + 1 + D
    counter = itertools.count()
    g = {(start, 0): 0}
    parent: Dict[Tuple[Cell, int], Tuple[Cell, int]] = {}
    open_heap = [(dist[start], 0, next(counter), (start, 0))]
    closed = set()
    while open_heap:
        f, neg_g, _, state = heapq.heappop(open_heap)
        if state in closed:
            continue
        closed.add(state)
        if stats is not None:
            stats.low_level_expanded += 1
        v, t = state
        if v == goal and t > goal_block:            # goal-stay test
            path = [v]
            while state in parent:
                state = parent[state]
                path.append(state[0])
            return path[::-1]
        if t >= horizon:
            continue
        for w in grid.neighbours(v) + [v]:           # moves and the wait action
            nxt = (w, t + 1)
            if nxt in closed or nxt in cons.vertex or (v, w, t) in cons.edge:
                continue
            if t + 1 < g.get(nxt, float("inf")):
                g[nxt] = t + 1
                parent[nxt] = state
                heapq.heappush(open_heap, (t + 1 + dist[w], -(t + 1), next(counter), nxt))
    return None
\end{lstlisting}

\begin{lstlisting}[caption={The body of the high level \code{cbs()}: root, best-first loop on the sum of costs, split, replan of one agent, optional bypass (from \code{code/ch09\_cbs.py}; \code{\_make\_child} unions the parent's constraints with the new one and calls \code{low\_level} for that agent only).},label={lst:ch09-highlevel}]
    t0, stats, trace, ids = time.perf_counter(), Stats(), [], itertools.count()
    dists = [grid.distances(g) for g in goals]         # one BFS per agent
    paths = [low_level(grid, s, g, Constraints(), d, stats)
             for s, g, d in zip(starts, goals, dists)]  # root: no constraints
    if any(p is None for p in paths):
        return Result(None, None, stats)
    root = CTNode(tuple(Constraints() for _ in starts), paths,
                  sum_of_costs(paths), len(all_conflicts(paths)), next(ids))
    open_heap = [(root.cost, root.n_conflicts, root.id, root)]
    stats.generated = 1
    while open_heap:
        _, _, _, node = heapq.heappop(open_heap)     # lowest cost, fewest conflicts
        stats.expanded += 1
        conflict, kids = _choose_conflict(node, split, grid, starts, goals, dists, stats)
        entry = _record(node, conflict, trace) if keep_trace else None
        if conflict is None:                          # goal node: conflict-free
            stats.seconds = time.perf_counter() - t0
            return Result(node.paths, node.cost, stats, trace)
        if (time_limit is not None and time.perf_counter() - t0 > time_limit) or \
           (node_limit is not None and stats.expanded >= node_limit):
            break
        children = []
        for kid, (agent, extra) in zip(kids, conflict.constraints()):
            if kid is None and split == "first":     # replan only that agent
                kid = _make_child(node, agent, extra, grid, starts, goals, dists, stats)
            if kid is not None:
                children.append(kid)
        if bypass and _bypass(node, children):
            stats.bypasses += 1
            if entry is not None:
                entry["bypass"] = True
            heapq.heappush(open_heap, (node.cost, node.n_conflicts, node.id, node))
            continue
        for kid in children:
            kid.id = next(ids)
            stats.generated += 1
            if entry is not None:
                entry["children"].append(_child_summary(kid))
            heapq.heappush(open_heap, (kid.cost, kid.n_conflicts, kid.id, kid))
    stats.seconds = time.perf_counter() - t0
    return Result(None, None, stats, trace)
\end{lstlisting}

\begin{pitfall}[Replan only the constrained agent, and do replan it]
The child of a split must (a)~inherit \emph{every} constraint of its parent,
(b)~add exactly one, and (c)~recompute the path of the constrained agent and
no other. Forgetting (a), for instance by building the child's set from the
new constraint alone, lets the agent walk back into a cell it was forbidden
earlier on the branch, and the tree loops through the same conflicts.
Forgetting (c) by copying the parent's path list by reference, or by
replanning the wrong agent of the pair, produces a node whose paths violate
its own constraints; the node's cost is then wrong, and \cref{thm:ch09-lowerbound}
fails, so the returned plan may be suboptimal or conflicting. Replanning
\emph{all} agents is not wrong but wastes $k-1$ searches per node and, with
low-level tie-breaking, can shuffle paths that had no reason to change.
Immutable constraint sets and an independent \code{validate} of every
returned plan make all three mistakes visible at once.
\end{pitfall}

\begin{pitfall}[A parked drone still occupies its goal]
Two places must know that an agent stays at its goal for ever. The conflict
check must compare agents beyond the end of the shorter path, treating the
arrived agent as sitting on its goal cell; otherwise a drone flying through
another's goal after its arrival is never reported and the plan is invalid.
And the low level must apply the goal-stay test
(line~\ref{alg:ch09-lowlevel:goal}): a constraint $\langle a_i, \gamma_i, t\rangle$
with $t$ later than the agent's arrival is only satisfied by a path that
arrives after $t$ or leaves and comes back. A low level that accepts the
first arrival returns the parent's path unchanged, the child has the same
conflict as the parent, and the high level splits the same conflict for
ever.
\end{pitfall}

\begin{pitfall}[An unbounded time horizon]
Space-time has infinitely many states, because waiting is always possible.
If the constraints seal an agent in, a low level without a horizon runs for
ever, expanding $(v, t)$ for larger and larger $t$; the same happens when a
constraint forbids the goal at a time that no path can avoid. The horizon
$T_{\max} = H + 1 + D$ of line~\ref{alg:ch09-lowlevel:horizon} keeps every
low-level call finite without excluding any optimal constrained path. The
high level has the corresponding problem on unsolvable instances
(\cref{rem:ch09-unsolvable}); always run it with a time or node limit.
\end{pitfall}

% ---------------------------------------------------------------------
\section{Where this fits in the drone system}
\label{sec:ch09-drone}
\begin{dronebox}
\cbs is the global swarm planner, the first of the four layers of the
hybrid architecture in \cref{ch:ch24}. Before take-off it receives the
inflated occupancy grid of \cref{ch:ch02}, one start and one goal cell per
drone, and a time step chosen so that a drone crosses one cell per step at
cruise speed, and it returns the nominal time-indexed routes
$\pi_1, \dots, \pi_k$, conflict-free and of minimum total flight time. In
three dimensions nothing changes except the successor function of the low
level (a 6-connected grid of altitude layers). The routes are what the
other layers work with: the prediction layer (\cref{ch:ch18,ch:ch20})
watches for intruders that were not in the plan, the local layer
(\cref{ch:ch13,ch:ch14}) lets a drone deviate from its route to avoid them,
and the replanning layer (\cref{ch:ch05}) repairs the deviated drone's
route, using the routes of the other drones as constraints in exactly the
form of the low level here. After any such repair the architecture runs the
conflict check of line~\ref{alg:ch09-cbs:validate} on the updated plan: if
the repaired route conflicts with another drone, the affected pair is
re-planned with \cbs on the constraints in force, and only then is the new
plan released. The optimality guarantee of \cref{thm:ch09-optimal} holds
for the nominal plan; after deviations, the re-check is what preserves
conflict-freedom. When more than a dozen drones interact, the global planner
is \ecbs (\cref{ch:ch10}) with the same two levels. This chapter is Week~4
of the study plan (\cref{ch:appA}): its milestone is the solver of
\cref{exr:ch09-coding}.
\end{dronebox}

% ---------------------------------------------------------------------
\section{Summary}
\begin{summary}
\begin{itemize}
  \item \cbs solves MAPF optimally with two levels: the high level searches a binary \emph{constraint tree} best-first on the sum of costs, the low level plans one agent with space-time \astar under that agent's constraints.
  \item A node holds constraint sets, one minimum-cost path per agent that respects them, and the sum of the path costs. The root has no constraints. A node whose paths are conflict-free is a goal node.
  \item A conflict $\langle a_i, a_j, v, t\rangle$ is split into two children, one adding $\langle a_i, v, t\rangle$ and one adding $\langle a_j, v, t\rangle$; an edge conflict (the swapping conflict of \cref{def:ch07-conflicts}) is split into the two opposite edge constraints. Only the constrained agent is replanned; constraints accumulate down the tree.
  \item Every solution satisfies at least one of the two constraints of a split (\cref{thm:ch09-split}), and a node's cost lower-bounds every solution below it (\cref{thm:ch09-lowerbound}); hence the first conflict-free node popped is optimal (\cref{thm:ch09-optimal}) and \cbs terminates on solvable instances (\cref{thm:ch09-complete}).
  \item The effort is exponential in the number of conflicts that must be resolved, not directly in the number of agents: \cbs excels at bottlenecks and corridors with few alternatives and suffers in open rooms and crowded mazes, where joint-space search can be better.
  \item The low level must obey the constraints \emph{and} the goal-stay test of \cref{ch:ch08}, and needs the horizon $T_{\max} = H + 1 + D$; a parked agent still conflicts.
  \item ICBS splits cardinal conflicts first (identified with MDDs) and bypasses non-cardinal ones; MA-CBS merges agents that conflict repeatedly; CBSH adds admissible high-level heuristics such as the conflict-graph vertex cover; rectangle and corridor reasoning replace one-cell constraints by barrier and range constraints. All keep optimality because they are legal splits.
  \item In the drone system \cbs computes the nominal time-indexed routes before flight, and its conflict check re-validates every repaired route.
\end{itemize}
\end{summary}

\paragraph{Further reading.}
The original conference paper~\cite{sharon2012cbs} and the journal
version~\cite{sharon2015cbs} contain the algorithm, the proofs, MA-CBS and
the corridor-versus-room analysis. ICBS is due to Boyarski et
al.~\cite{boyarski2015icbs}, the MDD to the increasing cost tree
search~\cite{sharon2013icts}, high-level heuristics to Felner et
al.~\cite{felner2018adding} and Li et al.~\cite{li2019improved}, and
symmetry reasoning to Li et al.~\cite{li2019symmetry,li2020pairwise}. The
MAPF definitions and benchmark instances follow Stern et
al.~\cite{stern2019mapf}; the joint-space baseline with operator
decomposition and independence detection is Standley's~\cite{standley2010finding}.
The bounded-suboptimal variants start with Barer et al.~\cite{barer2014ecbs}
and are the subject of \cref{ch:ch10}.

% ---------------------------------------------------------------------
\section{Exercises}
\label{sec:ch09-exercises}

\begin{exercise}[\difficulty{1}]
\label{exr:ch09-constraints}
In node $N_2$ of \cref{ex:ch09-crossing} the paths are $\pi_1 =
(0,1),(1,1),(2,1),(3,1),(4,1)$, $\pi_2 = (2,0),(2,1),(2,2)$ and $\pi_3 =
(2,2),(2,2),(2,1),(3,1),(4,1),(4,0)$. List all four conflicts of this node
(type, agents, cells, time) and, for each, the two constraints that
\cref{eq:ch09-split} would generate. Which of them does \cref{alg:ch09-cbs}
split, and why?
\end{exercise}

\begin{exercise}[\difficulty{2}]
\label{exr:ch09-handtree}
Consider the $3 \times 3$ grid in which the cells $(0,1)$ and $(2,1)$ are
blocked, so that the only way from the bottom row to the top row is through
the gap $(1,1)$. Agent $a_1$ flies from $(0,0)$ to $(2,2)$ and $a_2$ from
$(2,0)$ to $(0,2)$. Build the constraint tree by hand: give the root paths
and cost, the first conflict, the two children with their paths and costs,
and the returned plan. How does the size of this tree compare with the
number of joint states that \astar in the joint space would have to
consider?
\end{exercise}

\begin{exercise}[\difficulty{1}]
\label{exr:ch09-goalstay}
In the map of \cref{ex:ch09-crossing} let $a_1$ fly from $(0,1)$ to
$(4,1)$ as before, and let a second agent $a_2$ fly from $(0,0)$ to $(4,0)$,
which forces it through the whole horizontal corridor and through $a_1$'s
goal cell $(4,1)$ at $t = 5$, after $a_1$ has arrived at $t = 4$. Explain why
the root has a conflict, what constraint each child receives, and what the
low level does with the constraint $\langle a_1, (4,1), 5\rangle$ in
particular. What would happen if the low level accepted the first arrival
at the goal?
\end{exercise}

\begin{exercise}[\difficulty{2}]
\label{exr:ch09-proof-lb}
(a) \Cref{alg:ch09-cbs} splits the earliest conflict. Show, using only
\cref{thm:ch09-split,thm:ch09-lowerbound}, that \cbs stays optimal and
complete if it splits \emph{any} conflict of the node, for instance the
last one or a random one. (b) Show that discarding a child whose low-level
search fails (line~\ref{alg:ch09-cbs:replan}) never discards a solution.
(c) Give an instance in which the choice of the conflict changes the number
of expanded nodes.
\end{exercise}

\begin{exercise}[\difficulty{2}]
\label{exr:ch09-cardinal}
Verify, from the child costs listed in \cref{tab:ch09-ct}, that every
conflict split in \cref{ex:ch09-crossing} is cardinal. Then draw
the MDD of $a_1$ at the root and explain why the conflict
$\langle a_1, a_2, (2,1), 2\rangle$ of $N_1$ is cardinal although $a_1$'s
corridor is three cells wide at $x = 2$. Finally compute
$h_{\mathrm{CG}}$ of \cref{eq:ch09-hcg} for $N_0$, $N_1$ and $N_2$; for
$N_2$ you may assume that its three $a_1$--$a_3$ conflicts are cardinal,
so that you do not have to run six low-level searches by hand.
\end{exercise}

\begin{exercise}[\difficulty{2}]
\label{exr:ch09-rectangle}
In \cref{fig:ch09-symmetry} (left), $a_1$ flies from $(0,2)$ to $(4,3)$ and
$a_2$ from $(1,1)$ to $(3,4)$ on an obstacle-free $5 \times 5$ grid. Count
the shortest paths of each agent, and prove that every pair of shortest
paths has a vertex conflict. Hint: at time $t$ both agents are on the
anti-diagonal $x + y = 2 + t$; compare the $x$-coordinates. Conclude a lower
bound on the optimal sum of costs.
\end{exercise}

\begin{exercise}[\difficulty{2}]
\label{exr:ch09-makespan}
Suppose the high level orders the constraint tree by the makespan
$\max_i \cost(N.\pi_i)$ instead of the sum of costs, with the low level
unchanged. Is the returned plan makespan-optimal? Prove your answer by
adapting \cref{thm:ch09-lowerbound,thm:ch09-optimal}, and name one
practical difference from the sum-of-costs version.
\end{exercise}

\begin{exercise}[\difficulty{2}]
\label{exr:ch09-memory}
A constraint-tree node in \code{ch09\_cbs.py} stores $k$ paths of up to $T$
cells and $k$ constraint sets. Estimate the memory of a tree with $10^5$
generated nodes for $k = 20$ and $T = 60$. Then design a node that stores
only the replanned path, the added constraint and a parent pointer, and
give the cost of reconstructing the $k$ paths and the constraint sets of a
node at depth $d$.
\end{exercise}

\begin{exercise}[\difficulty{3} Coding]
\label{exr:ch09-coding}
The Week~4 exercise of the study plan. Implement \cbs from scratch for a
4-connected grid: a low level with vertex and edge constraints, the
goal-stay test and the horizon; conflict detection with parked
agents; and the high level with a heap keyed on the sum of costs. Reproduce
the tree of \cref{tab:ch09-ct}. Then test it on random $8 \times 8$ grids
with $2$, $4$ and $8$ agents, report the number of expanded nodes and the
running time, and certify optimality on tiny instances against a
brute-force search over the joint state space (as \code{ch09\_cbs.py}
does). Validate every returned plan independently.
\end{exercise}

\begin{exercise}[\difficulty{3} Coding]
\label{exr:ch09-icbs-coding}
Extend your solver from \cref{exr:ch09-coding} with cardinal-first
splitting and the bypass, first by generating both children of each
candidate conflict as \code{ch09\_cbs.py} does, then with MDDs built by two
breadth-first searches in space-time. Reproduce \cref{fig:ch09-scaling},
add a curve for the MDD-based classifier, and report how the number of
low-level expansions per constraint-tree node changes.
\end{exercise}
